% Generated draft: requires mathematical and source review.
\documentclass[11pt]{article}
\usepackage[margin=25mm]{geometry}
\usepackage{fontspec}
\setmainfont{Latin Modern Roman}
\usepackage{amsmath,amssymb,mathtools,unicode-math}
\setmathfont{Latin Modern Math}
\usepackage{enumitem,needspace,xurl,hyperref}
\setlength{\parindent}{0pt}
\setlength{\parskip}{5pt}
\setlength{\emergencystretch}{3em}
\newcommand{\N}{\mathbb{N}}
\newcommand{\Z}{\mathbb{Z}}
\newcommand{\Q}{\mathbb{Q}}
\newcommand{\R}{\mathbb{R}}
\newcommand{\C}{\mathbb{C}}
\newcommand{\F}{\mathbb{F}}
\newcommand{\A}{\mathbb{A}}
\newcommand{\PP}{\mathbb{P}}
\newcommand{\E}{\mathbb{E}}
\newcommand{\eps}{\varepsilon}
\newcommand{\dd}{\,\mathrm{d}}
\begin{document}
\section*{UnsolvedMath ID: 801130}
\textbf{Source problem code:} catalogue entry
\subsection{Definitions and mathematical statement}
\textbf{Context and definitions.} The source concerns ultraproducts of type $\mathrm{II}_1$ factors, but does not explicitly specify $N$, its trace, or $\mathcal U$. Under the standard contextual interpretation, let $N$ be a type $\mathrm{II}_1$ factor with normalized faithful normal trace $\tau$, and let $\mathcal U$ be a nonprincipal ultrafilter on $\mathbb N$. Its tracial ultrapower is
\[
N^{\mathcal U}=\ell^\infty(\mathbb N,N)\big/\mathcal I_{\mathcal U},\qquad
\mathcal I_{\mathcal U}=\left\{(x_n):\lim_{n\to\mathcal U}\|x_n\|_{2,\tau}=0\right\},
\]
where $\|x\|_{2,\tau}=\tau(x^*x)^{1/2}$ and $\tau_{\mathcal U}([(x_n)])=\lim_{n\to\mathcal U}\tau(x_n)$. These contextual hypotheses require confirmation from the source.

Unitaries $u,v\in N^{\mathcal U}$ are freely independent if their generated unital $*$-algebras $A_u$ and $A_v$ satisfy
\[
\tau_{\mathcal U}(a_1\cdots a_m)=0
\]
for every integer $m\geq 1$ and every choice of centered elements $a_j\in A_{\epsilon_j}$ with $\epsilon_j\in\{u,v\}$ and $\epsilon_j\ne\epsilon_{j+1}$. Here centered means $\tau_{\mathcal U}(a_j)=0$.

\textbf{Source question.} Suppose $u,v\in N^{\mathcal U}$ are freely independent unitaries. Can one find sequences of unitaries $(u_n)$ and $(v_n)$ in $N$ such that
\[
u=[(u_n)]_{\mathcal U},\qquad v=[(v_n)]_{\mathcal U},
\]
and the sequences have the source's stated property of being ``pairwise $k$-independent unitaries in $N$ for all $k$''?

\textbf{Unresolved formulation.} The supplied record does not define $k$-independence or the meaning of ``pairwise'' for these sequences. It also does not specify the domain of $k$, whether one pair of representing sequences must work simultaneously for every $k$, or whether the required independence holds at every coordinate, eventually, or on a set belonging to $\mathcal U$. Consequently the quoted independence requirement cannot be replaced by a precise formula from the supplied evidence alone.

The source says that particular cases are established in the reference labeled \texttt{[exoticCIKE]}. Its contextual note attributes a solution to G. Patchell and A. Shiner, with a paper in preparation, whereas the registry status is \texttt{open}. These are supplied source qualifications, not independently verified mathematical conclusions.
\subsection{Short English statement}
The question asks whether freely independent unitaries in a tracial ultrapower can be represented by unitary sequences satisfying the source's pairwise $k$-independence requirement for all $k$. The meaning and quantifiers of that requirement are unspecified. The registry labels the problem open, while its contextual note reports a solution by G. Patchell and A. Shiner in a paper in preparation.
\subsection{Sources}
\begin{itemize}
\item Supplied local UnsolvedMath registry record: permanent ID $801130$, exact source problem code \texttt{catalogue entry}; upstream revision \texttt{d371ee917ac9647fbb06cd6fc529deec6f3dda82}. The record attributes the problem to Elayavalli, with proposed year $2025$.
\item Elayavalli, \emph{50 Open Problems: Ultraproduct II\_1 Factors} ($2025$), Problem $31$: \url{https://arxiv.org/abs/2511.20377}. The supplied provenance records access on 2026-10-16 and extraction by \texttt{arxiv\_latex}; these provenance claims have not been independently verified here.
\item The supplied statement cites \texttt{[exoticCIKE]} for particular cases. No bibliographic identification or link for this citation is supplied.
\item Supplied attribution/context: ``Solved by G. Patchell and A. Shiner, paper in preparation.'' No title, manuscript link, or proof reference for that reported solution is supplied. The registry separately records status \texttt{open}.
\end{itemize}
\end{document}
